\section{Model and extremal parameters}\label{sec:general-model}\label{sec:model}

Identify the compass directions with the unit vectors
\[
 \N=(0,1),\qquad \E=(1,0),\qquad \Sdir=(0,-1),\qquad \W=(-1,0),
 \qquad D=\{\N,\E,\Sdir,\W\}.
\]
Juxtaposition abbreviates a mask, for example $\N\Sdir=\{\N,\Sdir\}$.
A \emph{maze} is a finite nonempty set $V\subset\Z^2$ whose induced grid graph is connected.  Thus every pair of cells of Manhattan distance one is joined by an edge.  The mask observed at $x\in V$ is
\[
 M_V(x)=\{d\in D:x+d\in V\}.
\]
We write $d_V$ for graph distance in the induced grid graph and $d_1(x,y)=\|x-y\|_1$ for Manhattan distance; always $d_1(x,y)\le d_V(x,y)$.

\begin{definition}[Controller and execution]\label{def:controller-general}
A controller $A$ consists of a finite nonempty state set $Q_A$, a distinguished initial state $0$, and a deterministic legal transition table
\[
 \delta_A:Q_A\times(2^D\setminus\{\varnothing\})\longrightarrow D\times Q_A,
 \qquad \delta_A(q,M)=(d,r),\quad d\in M.
\]
The direction is written first and the new state second.  Both agents use the same table and the same global compass, start in state $0$, and move synchronously.  There is no wait action, marking, communication, external clock, or access to coordinates.  For $|V|\ge2$ the one-agent map is
\[
 \Phi_{A,V}(x,q)=(x+d,r),\qquad (d,r)=\delta_A(q,M_V(x)).
\]
Starting from cells $x,y\in V$, the two copies iterate this map in parallel from $(x,0)$ and $(y,0)$.
\end{definition}

Fix a real success radius $\rho\ge0$.  The execution is successful if, for some integer time $t\ge0$,
\[
 d_1(X_t,Y_t)\le\rho.
\]
A pair that is already within radius $\rho$ succeeds at time zero.  Passing through the same edge in opposite directions creates no additional continuous-time event.  Since Manhattan distances on the grid are integral,
\begin{equation}\label{eq:integer-radius}
 d_1(X_t,Y_t)\le\rho
 \quad\Longleftrightarrow\quad
 d_1(X_t,Y_t)\le\lfloor\rho\rfloor .
\end{equation}
A singleton has only a time-zero success and requires no transition at the empty mask.
Thus all three parameters below are unchanged when $\rho$ is replaced by $\lfloor\rho\rfloor$.

Let $\mathcal G$ denote all finite connected induced grid mazes, $\mathcal T_{\rm grid}$ all induced grid trees, $\mathcal P$ all induced grid paths, and $\mathcal J$ the class consisting of paths together with induced grid trees having exactly one degree-three vertex and all other degrees at most two.  Then
\[
 \mathcal P\subseteq\mathcal J\subseteq\mathcal T_{\rm grid}\subseteq\mathcal G.
\]

\begin{definition}[Controller-dependent trap size and $F$]\label{def:F-general}
For a nonempty maze class $\mathcal C$, let
\[
 \tau_{\mathcal C,\rho}(A)
 =\min\Bigl\{|V|:\ V\in\mathcal C,\ \exists x,y\in V\ 
       \forall t\ge0,\ d_1(X_t,Y_t)>\rho\Bigr\},
\]
with $\min\varnothing=\infty$.  Thus $\tau_{\mathcal C,\rho}(A)=\infty$ means that the one fixed controller $A$ succeeds from every start pair on every maze in $\mathcal C$.  For a state budget $s\ge1$, define
\[
 \boxed{F_{\mathcal C,\rho}(s)
 =\max_{1\le |Q_A|\le s}\tau_{\mathcal C,\rho}(A).}
\]
The controller is chosen first; the adversary then chooses the maze and the starts.  The trapping maze may therefore depend on the controller.
\end{definition}

For a fixed state budget there are finitely many labelled legal transition tables.  Hence the maximum in Definition~\ref{def:F-general} is attained in $\mathbb N\cup\{\infty\}$.  In particular,
\begin{equation}\label{eq:F-infinity-meaning}
 F_{\mathcal C,\rho}(s)=\infty
 \quad\Longleftrightarrow\quad
 \text{one controller with at most $s$ states succeeds on all of $\mathcal C$.}
\end{equation}
The value $\infty$ is therefore not an unattained supremum of finite trap sizes.

\begin{definition}[Finite-horizon state complexity]\label{def:M-general}
For $n\in\mathbb N_0$, let
\[
 \boxed{M_{\mathcal C,\rho}(n)
 =\min\Bigl\{|Q_A|:\ \forall V\in\mathcal C,\ |V|\le n,\ 
       \forall x,y\in V,\ \exists t\ge0:\ d_1(X_t,Y_t)\le\rho\Bigr\},}
\]
again with an empty minimum equal to $\infty$.  Thus $M_{\mathcal C,\rho}(n)$ is the least number of states in one controller that succeeds on every legal instance through size $n$.
\end{definition}

The strict threshold relation between the two parameters is immediate but useful enough to record at once.

\begin{theorem}[$F/M$ duality]\label{thm:FM-duality}
For every finite $n$ and every finite state budget $s\ge1$,
\[
 \boxed{F_{\mathcal C,\rho}(s)>n
 \quad\Longleftrightarrow\quad
 M_{\mathcal C,\rho}(n)\le s.}
\]
\end{theorem}
\begin{proof}
By Definition~\ref{def:F-general}, $F_{\mathcal C,\rho}(s)>n$ holds exactly when there is a controller $A$ with $|Q_A|\le s$ and $\tau_{\mathcal C,\rho}(A)>n$.  The latter inequality says that no maze in $\mathcal C$ of size at most $n$, with any start pair, is a trap for $A$.  Equivalently,
\[
 \exists A,\ |Q_A|\le s:\ \forall V\in\mathcal C,\ |V|\le n,\ 
 \forall x,y\in V,\ \exists t\ge0:\ d_1(X_t,Y_t)\le\rho,
\]
which is precisely $M_{\mathcal C,\rho}(n)\le s$.  The equivalence also covers $M=\infty$ and $F=\infty$ under the conventions above.
\end{proof}

\begin{proposition}[Monotonicities]\label{prop:general-monotonicities}
The parameters are monotone in the directions shown in Table~\ref{tab:monotonicities}.  In addition,
\[
 F_{\mathcal C,\rho}(s)=F_{\mathcal C,\lfloor\rho\rfloor}(s),
 \qquad
 M_{\mathcal C,\rho}(n)=M_{\mathcal C,\lfloor\rho\rfloor}(n).
\]
\end{proposition}
\begin{proof}
Increasing the state budget enlarges the controller set over which $F$ is maximized.  Increasing $\rho$ can only turn failures into successes, so for each fixed controller its least trap cannot become smaller; maximizing preserves this direction.  If $\mathcal C_1\subseteq\mathcal C_2$, the adversary has fewer candidate traps in $\mathcal C_1$, so $\tau_{\mathcal C_1,\rho}(A)\ge\tau_{\mathcal C_2,\rho}(A)$ for every $A$, and hence the same inequality holds for $F$.

For $M$, enlarging the size horizon or the maze class strengthens the universal success requirement and therefore cannot decrease the required number of states; enlarging the success radius weakens that requirement and therefore cannot increase it.  The floor identities are exactly \eqref{eq:integer-radius}.
\end{proof}

\begin{table}[!htbp]
\centering\small
\begin{tabular}{@{}lll@{}}
\toprule
Change & $F$ & $M$ \\
\midrule
$s_1\le s_2$ & $F_{\mathcal C,\rho}(s_1)\le F_{\mathcal C,\rho}(s_2)$ & --- \\
$\rho_1\le\rho_2$ & $F_{\mathcal C,\rho_1}(s)\le F_{\mathcal C,\rho_2}(s)$ & $M_{\mathcal C,\rho_1}(n)\ge M_{\mathcal C,\rho_2}(n)$ \\
$\mathcal C_1\subseteq\mathcal C_2$ & $F_{\mathcal C_1,\rho}(s)\ge F_{\mathcal C_2,\rho}(s)$ & $M_{\mathcal C_1,\rho}(n)\le M_{\mathcal C_2,\rho}(n)$ \\
$n_1\le n_2$ & --- & $M_{\mathcal C,\rho}(n_1)\le M_{\mathcal C,\rho}(n_2)$ \\
\bottomrule
\end{tabular}
\caption{Monotonicities.  Larger state budgets and success radii help the controller, whereas enlarging the maze class helps the adversary.  For $M$, increasing the size horizon makes the requirement harder.}
\label{tab:monotonicities}
\end{table}


For the original model we abbreviate
\[
 F(s)=F_{\mathcal G,1}(s),\qquad M(n)=M_{\mathcal G,1}(n).
\]
At radius one we also use the fixed-radius shorthands
$\tau_{\mathcal C}(A)=\tau_{\mathcal C,1}(A)$ and
$F_{\mathcal C}(s)=F_{\mathcal C,1}(s)$, and write $\tau(A)=\tau_{\mathcal G,1}(A)$.
For a fixed controller, $\Phi_V$ abbreviates $\Phi_{A,V}$ and $x_q$ denotes the tag $(x,q)$.
An initially noncontacting start-labelled maze whose execution never succeeds is a
\emph{trap instance}. These are shorthands for the same model, not new parameters.
\begin{equation}\label{eq:class-monotonicity}
\mathcal C_1\subseteq\mathcal C_2\quad\Longrightarrow\quad
\tau_{\mathcal C_1,\rho}(A)\ge\tau_{\mathcal C_2,\rho}(A),\qquad
F_{\mathcal C_1,\rho}(s)\ge F_{\mathcal C_2,\rho}(s).
\end{equation}

\begin{lemma}[Simultaneous compass transport]\label{m:transport}
Let $g\in D_4$ be a rotation or reflection of the square grid. Define the transported controller $gA$ by
\[
 \delta_{gA}(q,gM)=(gd,r)\quad\hbox{when}\quad\delta_A(q,M)=(d,r).
\]
Transporting the maze and the starts by $g$ carries every execution of $A$ to an execution of $gA$ with the same internal states and the same contact times.
\end{lemma}
\begin{proof}
The induced mask at $gx$ in $gV$ is $gM_V(x)$. Thus $\Phi_{gV}^{gA}(gx,q)=(g(x+d),r)$ whenever $\Phi_V^A(x,q)=(x+d,r)$. Iteration proves the execution identity, and $g$ preserves Manhattan distance.
\end{proof}

The lemma permits compass normalization of a controller-dependent upper witness. It does not permit rotating a maze while leaving a fixed controller unchanged. In the lower proof every orientation is absolute, and the initial state is never exchanged with state~$1$.
